The configured supervisor stack pairs select only the first user-origin entry and delivery-failure recovery.

\manualtablecaption{Supervisor Stack Roles}
\begin{manualtabular}{@{}>{\raggedright\arraybackslash}p{1.10in}>{\raggedright\arraybackslash}p{1.15in}>{\raggedright\arraybackslash}p{3.15in}@{}}
\toprule
\rowcolor{ManualHeaderFill}
\textbf{Pair} & \textbf{Role} & \textbf{Use}\\
\midrule
\texttt{SSS:SSP} & system call & first user-origin SYSTEM\_CALL\\
\texttt{ISS:ISP} & interrupt & first user-origin maskable interrupt\\
\texttt{FSS:FSP} & fault & first user-origin exception or NMI\\
\texttt{DSS:DSP} & recovery & DOUBLE\_FAULT after delivery failure\\
\bottomrule
\end{manualtabular}

A first user-origin entry loads the applicable configured pair. Supervisor-origin and all nested entries retain the current SS:SP and push a full frame there; event class never promotes an active handler to another configured stack. Delivery failure uses DSS:DSP when the hidden current-DFA bit is clear. Failure with current DFA set enters shutdown.

For user-origin entry, hardware rounds only the event payload to 16 bytes and subtracts it from the configured top. BASIC events have zero payload and perform no stack access. For supervisor or nested entry, hardware subtracts the complete full-frame size from current SP. Each nonzero range is validated and stored atomically before the new SP commits; configured tops are never modified.
